\ExerciseSection

\begin{enumerate}
\def\labelenumi{\arabic{enumi})}
\item
  设 \(f\) 是 \(\mathbf{R}\) 上的实值函数，它等于 \(0\) （当 \(x \leqslant 0\) ），等于 \(x\) （当 \(0 \leqslant x \leqslant 1\) ），等于 \(1\) （当 \(x \geqslant 1\) ）。证明实值函数序列 \(f_n(x) = f(nx - n^2)\) ( \(n \in \mathbf{N}\) ) 是等度连续的但在 \(\mathbf{R}\) 上不是一致等度连续的，尽管它由一致连续函数组成。
\item
  设 X 是每个点都具有可数基本邻域系的豪斯多夫拓扑空间，Y 是一致空间。设 H 是 \(\mathcal{C}(X;Y)\) 的一个子集，具有下列性质：对 X 的每个紧子集 K，集 \(H|K\) (由映射 \(u\in H\) 在 K 上的限制组成)是 \(\mathcal{C}(K;Y)\) 的一个等度连续子集。证明 H 是等度连续的（用反证法论证）。
\item
  设 X、Y、Z 是三个度量空间，H 是 \(\mathcal{C}(X \times Y; Z)\) 的一个子集。假设对每个 \(x_{0} \in X\) ，当 u 跑遍 H 时映射 \(u(x_{0}, .)\) 所成之集是 \(\mathcal{C}(Y; Z)\) 的一个等度连续子集，且对每个 \(y_{0} \in Y\) ，映射 \(u(., y_{0})\) 所成之集当 u 跑遍
\end{enumerate}

H 时是 \(\mathcal{C}\) (X; Z) 的一个等度连续子集。证明若 X 是完备的，则对每个 \(b \in \mathbf{Y}\) ，存在 X 中的一个集 \(S_b\) ，它的余集在 X 中是贫集，且对每个 \(a \in \mathbf{S}_b\) ，集 H 在 \((a, b)\) 处等度连续。（应用第 IX 章，§ 5，习题 23，并利用第 X 章，§ 2，命题 1，推论 1。）

\begin{enumerate}
\def\labelenumi{\arabic{enumi})}
\setcounter{enumi}{3}
\tightlist
\item
  设 X、Y、Z 是三个一致空间。
\end{enumerate}

\begin{enumerate}
\def\labelenumi{\alph{enumi})}
\item
  设 H 是 \(\mathcal{C}(\mathrm{Y};\mathrm{Z})\) 的一个一致等度连续子集。若 H、\(\mathcal{C}(\mathrm{X};\mathrm{Y})\) 与 \(\mathcal{C}(\mathrm{X};\mathrm{Z})\) 都赋有一致收敛的一致结构，证明映射 \((u,v)\to u\circ v\) (从 \(\mathrm{H}\times \mathcal{C}(\mathrm{X};\mathrm{Y})\) 到 \(\mathcal{C}(\mathrm{X};\mathrm{Z})\) 中)是一致连续的。
\item
  设 K 是 \(\mathcal{C}(\mathbf{X};\mathbf{Y})\) 的一个一致等度连续子集，L 是 \(\mathcal{C}(\mathbf{Y};\mathbf{Z})\) 的一个一致等度连续子集。证明当 u 跑遍 K 而 v 跑遍 L 时映射 \(v \circ u\) 所成之集是 \(\mathcal{C}(\mathbf{X};\mathbf{Z})\) 的一个一致等度连续子集。
\end{enumerate}

\begin{enumerate}
\def\labelenumi{\arabic{enumi})}
\setcounter{enumi}{4}
\item
  设 X 是拓扑空间，Y 是非离散赋值除环 K 上的赋范空间。设 H 是 \(\mathcal{F}(X;Y)\) 的一个在一点 \(x_{0}\in X\) 处等度连续的子集，并设 k\textgreater0 是一个实数。则集 \(H_{k}\) (由线性组合 \(\sum_{i}c_{i}u_{i}\) 组成，它们是函数 \(u_{i}\in H\) 的满足 \(\sum_{i}|c_{i}|\leqslant k\) 的线性组合)在点 \(x_{0}\) 处等度连续。
\item
  设 X 是拓扑空间，Y 是一致空间，H 是 \(\mathcal{C}(\mathbf{X};\mathbf{Y})\) 的一个等度连续子集，\(\Phi\) 是 H 上的一个滤子。证明 \(x\in X\) 中使 \(\Phi(x)\) 是 Y 上的柯西滤子基的那些 x 全体所成之集在 X 中是闭的。
\item
  设 X 是拓扑空间，Y 是完备豪斯多夫一致空间，\(\varphi\) 是 Y 到一个开子集 \(\varphi(Y)\) (它是一个豪斯多夫一致空间 \(Y'\) 的开子集)上的一个一致连续同胚。设 H 是 \(\mathcal{C}(X; Y)\) 的一个等度连续子集，\(H'\) 是映射 \(\varphi \circ u\) ( \(u \in H\) )全体所成之集。若 \(v\) 属于 \(H'\) 在 \(\mathcal{F}_s(X; Y')\) 中的闭包，证明 \(\overline{v}^{1}(\varphi(Y))\) 在 X 中既开又闭。特别地，若 X 连通，则 \(v(X)\) 含于 \(\varphi(Y)\) 或含于 \(\mathbb{C}\varphi(Y)\) 的一个连通分支。（注意 \(v\) 是连续的，并利用习题 6。）
\item
  设 \(\mathbf{H}\) 是拓扑空间 \(\mathbf{X}\) 到 \(\mathbf{R}\) 中的等度连续映射所成的一个集。
\end{enumerate}

\begin{enumerate}
\def\labelenumi{\alph{enumi})}
\item
  证明 H 的有限子集的上（相应地，下）包络所成之集是等度连续的。
\item
  设 v 是 X 到 \(\overline{R}\) 中的一个映射，它属于 H 关于逐点收敛拓扑的闭包。证明 v 在 \(\mathbf{X}\) 上连续，且集 \(\overline{v} (+\infty)\) 与 \(\overline{v} (-\infty)\) 在 \(\mathbf{X}\) 中都既开又闭（利用习题 7）。
\item
  由 a) 与 b) 推出，H 的上（相应地，下）包络 w（相应地，v）在 X 上连续，且 \(\overline{w}^{1}(+\infty)\) {[}相应地 \(\overline{v}^{1}(-\infty)\) {]}在 X 中既开又闭。
\item
  假设 X 连通且存在一点 \(x_{0} \in X\) 使 \(\mathbf{H}(x_{0})\) 在 R 中有上界。证明对 X 的每个紧子集 K，H 中函数在 K 上的限制所成之集在 K 中一致有上界{[}利用 c){]}。
\end{enumerate}

\begin{enumerate}
\def\labelenumi{\arabic{enumi})}
\setcounter{enumi}{8}
\item
  设 X 是一个集，Y 是一致空间，S 是 X 的一个覆盖。一致空间 \(\mathcal{F}_{\mathfrak{S}}(\mathbf{X};\mathbf{Y})\) 的子集 H 是预紧的，当且仅当 (i) 对每个 \(x\in X\) ，\(\mathbf{H}(x)\) 在 Y 中预紧；(ii) S-收敛的一致结构与逐点收敛的一致结构在 H 上一致。（把 X 看作离散空间并化归为 Y 是完备豪斯多夫的情形，应用定理 2。）
\item
  \begin{enumerate}
  \def\labelenumii{\alph{enumii})}
  \tightlist
  \item
    设 X 是紧空间，\((f_{n})\) 是连续映射 \(X \rightarrow X\) 的一个序列，它在 X 上逐点收敛但不一致收敛于一个连续函数（§ 1，第 6 号，注 2）。证明映射 \(f_{n}\) 所成之集在 \(\mathcal{C}_{s}(X; X)\) 中是相对紧的，但不相对紧于
  \end{enumerate}
\end{enumerate}

\[
\mathcal {C} _ {u} (\mathrm{X}; \mathrm{X}) = \mathcal {C} _ {c} (\mathrm{X}; \mathrm{X}).
\]

\begin{enumerate}
\def\labelenumi{\alph{enumi})}
\setcounter{enumi}{1}
\tightlist
\item
  设 I 是 R 的区间 \([-1, +1]\) ，且对每个整数 n \textgreater{} 0 ，令 \(u_{n}(x) = \sin \sqrt{x + 4n^{2}\pi^{2}}\) 对一切 \(x \geqslant 0\) 成立。证明函数 \(u_{n}\) 所成之集 H 是 \(\mathcal{C}(\mathbf{R}_{+}; \mathbf{I})\) 的一个等度连续子集，且在 \(\mathcal{C}_{c}(\mathbf{R}_{+}; \mathbf{I})\) 中相对紧但在 \(\mathcal{C}_{u}(\mathbf{R}_{+}; \mathbf{I})\) 中不相对紧。（注意序列 \((u_{n})\) 逐点收敛于 0。）
\end{enumerate}

\begin{enumerate}
\def\labelenumi{\arabic{enumi})}
\setcounter{enumi}{10}
\item
  设 X 是完全正则空间，S 是 X 的子集所成的一个集，这些子集的内部覆盖 X。证明空间 \(\mathcal{C}_{s}(X;R)\) 不是局部紧的。
\item
  设 X 是拓扑空间，Y 是一致空间，H 是 C(X; Y) 的一个等度连续子集，V 是 Y 的一个对称围域。给定一点 \(x \in X\) ，证明点 \(x' \in X\) (对其中每个点存在一个依赖于 \(x'\) 的整数 n 使得 \(\mathrm{H}(x') \subset \mathrm{V}^n(\mathrm{H}(x))\) )全体所成之集在 X 中既开又闭。由此推出，若 K 是 X 的任一紧连通子集且 \(x_0\) 是 K 的任一点，则存在整数 n \textgreater{} 0 使 \(\mathrm{H}(\mathrm{K}) \subset \mathrm{V}^n(\mathrm{H}(x_0))\) 。
\end{enumerate}

¶ 13) 设 X 是拓扑空间，Y 是局部紧一致空间，H 是 C(X;Y) 的一个等度连续子集。

\begin{enumerate}
\def\labelenumi{\alph{enumi})}
\item
  设 A 是 \(x \in X\) 中使 \(H(x)\) 在 Y 中相对紧的那些 x 全体所成之集。证明 A 在 X 中是开的（参见第 I 章，§ 9，第 7 号，命题 10 与第 II 章，§ 4，第 3 号，命题 4）。
\item
  再假设 Y 关于其一致结构是完备的；这时 A 在 X 中也是闭的{[}注意若 \(x_{0} \in \overline{A}\) ，则把 \(\mathrm{H}(x_{0})\) 上的一个超滤子看作 H 上一个超滤子的象，即可知 \(\mathrm{H}(x_{0})\) 在 Y 中相对紧{]}。在这一情形，若 X 连通，则 H 在 \(\mathcal{C}_{c}(\mathrm{X};\mathrm{Y})\) 中相对紧，当且仅当对一点 \(x_{0} \in X\) ，集 \(\mathrm{H}(x_{0})\) 在 Y 中相对紧。
\item
  取 X 为 R 的紧区间 {[}0, 1{]} ，取 Y 为赋有由 R 的一致结构诱导的一致结构的区间 {]}0, 1{[} 。给出 C (X; Y) 的一个等度连续子集 H 的例子，使得 a) 中定义的集 A 是区间 {]}0, 1{]} 。
\item
  假设 X = Y 且 H 由 X 到其自身上的同胚组成。证明若 H 是一致等度连续的，则 a) 中定义的集 A 在 X 中既开又闭。
\end{enumerate}

\begin{enumerate}
\def\labelenumi{\arabic{enumi})}
\setcounter{enumi}{13}
\tightlist
\item
  设 \(\Gamma\) 是 \(\mathbf{R}\) 的同胚所成的一个等度连续群。证明若同胚 \(u \in \Gamma\) 至少保持 \(\mathbf{R}\) 的一个点不动，且 \(u\) 是递增的，则 \(u\) 是恒等映射（证明否则由 \(u\) 生成的群在 \(u\) 的一个适当选取的不动点处不是等度连续的）。若 \(u\) 是递减的，则 \(u^2\) 是恒等映射。
\end{enumerate}

¶ 15) 设 X 是紧可度量化空间，\(\Gamma\) 是 X 的所有同胚所成的群，G 是 \(\Gamma\) 的一个在 X 上传递作用的子群。若 H 是 G 在 \(\Gamma\) 中的中心化子，证明 H 是等度连续的。{[}用归谬法论证，假设 H 在某点 \(a \in \mathbf{X}\) 处不是等度连续的；推出存在点的序列 \(x_n \in \mathbf{X}\) 与元素的序列 \(u_n \in \mathbf{H}\) 使得

\[
\lim _ {n \rightarrow \infty} x _ {n} = a, \quad \lim _ {n \rightarrow \infty} u _ {n} (a) = b, \quad \lim _ {n \rightarrow \infty} u _ {n} (x _ {n}) = c,
\]

其中 \(b \neq c\) 。由此证明序列 \((u_n)\) 在 \(X\) 中逐点收敛，但没有一个 \(X\) 的点是该序列的一致收敛点；这与 § 1 的习题 9 矛盾。{]}

\begin{enumerate}
\def\labelenumi{\arabic{enumi})}
\setcounter{enumi}{15}
\tightlist
\item
  设 X 是豪斯多夫一致空间，\(\Gamma\) 是 X 的同胚所成的一个等度连续群。由 \(\Gamma\) 在 X 上定义的等价关系 R 是开的（第 I 章，§ 5，第 2 号）。
\end{enumerate}

\begin{enumerate}
\def\labelenumi{\alph{enumi})}
\item
  证明若 \(\Gamma\) 的每个轨道在 \(\mathbf{X}\) 中是闭的，则轨道空间 \(\mathbf{X} / \Gamma\) 是豪斯多夫的。
\item
  证明若 \(\Gamma\) 的每个轨道是紧的，则关系 \(\mathbf{R}\) 是闭的（利用第 I 章，§ 5，第 4 号，命题 10）。
\item
  给出一个 X 紧但 \(\Gamma\) 的轨道没有一个在 X 中是闭的例子（参见第 III 章，§ 2，习题 29）。
\end{enumerate}

¶ 17) 设 X 是紧空间，Γ 是 X 的同胚所成的一个可数群。假设轨道空间 X/Γ 是豪斯多夫的（从而是紧的）。

\begin{enumerate}
\def\labelenumi{\alph{enumi})}
\item
  证明对每个 \(x \in \mathbf{X}\) ，轨道 \(\Gamma(x)\) (即 \(x\) 的轨道)是有限的（注意它若是无限的就没有孤立点，并利用贝尔定理（第 IX 章，§ 5，第 3 号，定理 1））。
\item
  对每个 \(x_0 \in \mathbf{X}\) ，设 \(\Delta(x_0)\) 是 \(\Gamma\) 的由使轨道 \(\Gamma(x_0)\) 的所有点都不动的同胚组成的正规子群。证明若 \(\Delta(x_0)\) 是有限生成的，则 \(\Gamma\) 在 \(x_0\) 处等度连续。{[}设 \(f_i (1 \leqslant i \leqslant m)\) 是 \(\Delta(x_0)\) 的生成元，并设 \(g_k (1 \leqslant k \leqslant n)\) 是 \(\Delta(x_0)\) 在 \(\Gamma\) 中除 \(\Delta(x_0)\) 自身以外的每个陪集的代表。取一个邻域 \(\mathbf{V}\) ( \(x_0\) 的邻域)使诸集 \(f_i(\mathbf{V})\) 、 \(f_i^{-1}(\mathbf{V})\) 中没有一个与诸集 \(g_k(\mathbf{V})\) 相遇，然后利用由 \(\Gamma\) 定义的等价关系是闭的这一事实（第 I 章，§ 10，第 4 号，命题 8）。{]}
\item
  不对 \(\Delta(x_{0})\) 作任何假设，而设 \(x_{0}\) 在 X 中有一个可数连通邻域基本系。证明 \(\Gamma\) 在 \(x_{0}\) 处等度连续{[}与 b) 同法{]}。推出若 X 是局部连通的，则 \(\Gamma\) 是等度连续的。
\item
  取 X 为 R 的由点 0, 1, 1/n 与 \(1 + 1/n\) (n 为整数 \(\geqslant 2\) )组成的紧子空间。给出 X 的同胚所成的一个可数群 \(\Gamma\) 的例子，使得 \(X/\Gamma\) 是豪斯多夫的但 \(\Gamma\) 不是等度连续的。
\end{enumerate}

¶ 18) 设 G 是连续作用在豪斯多夫拓扑空间 X 上的豪斯多夫拓扑群。称 G 在一点 \(x_{0} \in X\) 处是真作用的，如果轨道 G. \(x_{0}\) 在 X 中是闭的且 G 在 G. \(x_{0}\) 上真作用（第 III 章，§ 4，第 1 号）。

\begin{enumerate}
\def\labelenumi{\alph{enumi})}
\item
  设 G 是连续作用在豪斯多夫一致空间 X 上的局部紧拓扑群。假设当 s 跑遍 G 时 X 的同胚 \(x \rightarrow s.x\) 所成之集是等度连续的。证明若 G 在一点 \(x_{0} \in X\) 处真作用，则存在 \(x_{0}\) 的一个邻域 V，使得元素 \(s \in G\) (满足 \(s.V \cap V \neq \emptyset\) )全体所成之集在 G 中相对紧（参见第 III 章，§ 4，第 4 号，命题 7）。由此推出 X 中使 G 真作用的点所成之集 D 在 X 中是开的，且 G 在 D 上真作用。若此外 X 是局部紧的，证明 D 在 X 中既开又闭（利用第 1 号，命题 1，推论 2）。
\item
  在实平面 \(\mathbf{R}^2\) 中，设 \(\mathbf{E}\) 是由原点 \((0,0)\) 与点 \((0,2^{-n})\) (当 \(n\) 跑遍整数 \(\geqslant 0\) )组成的集。对每个 \(n\in \mathbf{Z}\) (满足 \(n\neq 0\) )，设 \(u_{n}\) 是一个仿射线性映射在 \(\mathbf{E}\) 上的限制 (该映射是 \(\mathbf{R}^2\) 到其自身中的映射)，使得 \(u_{n}(0,0) = (n,0)\) ，且 \(u_{n}(0,2^{-|n|}) = (0,\theta^{n})\) ，其中 \(\theta\) 是满足 \(0 <   \theta <  1\) 的一个无理数。设 \(\mathbf{X}\) 是 \(\mathbf{R}^2\) 的（局部紧）子空间，它是 \(\mathbf{E}\) 与诸集 \(u_{n}(\mathbf{E})\) ( \(n\in \mathbf{Z},n\neq 0\) )的并。此外，设 \(u_{0}\) 表示 \(\mathbf{E}\) 的恒等映射，并把 \(u_{n}(n\in \mathbf{Z})\) 延拓到整个 \(\mathbf{X}\) 上，其法是置 \(u_{n}(u_{m}(x)) = u_{n + m}(x)\) 对一切 \(x\in \mathbf{E}\) 与一切 \(m\in \mathbf{Z}\) 。集 \(\mathbf{G}\) (它由诸映射 \(u_{n}\) 组成)是 \(\mathbf{X}\) 的同胚所成的一个群，赋以离散拓扑使 \(\mathbf{G}\) 连续地作用在 \(\mathbf{X}\) 上。证明 \(\mathbf{G}\) 在 \(\mathbf{X}\) 的每一点处真作用，但不是等度连续的且不在 \(\mathbf{X}\) 上真作用。
\item
  设 X 是 \(R^{2}\) (等同于 C)的（非局部紧）子空间，它是半平面 y \textgreater{} 0 与原点的并。通过置 \(u(0, 0) = (0, 0)\) 与 \(u(re^{i\omega}) = re^{i\omega'}\) 定义 X 到其自身上的一个同胚 u，其中 \(\omega' = \frac{1}{2}\omega\) 当 \(0 < \omega \leqslant \frac{1}{2}\pi\) ，\(\omega' = \omega - \frac{1}{4}\pi\) 当 \(\frac{1}{2}\pi \leqslant \omega \leqslant \frac{3}{4}\pi\) ，\(\omega' = 2\omega - \pi\) 当 \(\frac{3}{4}\pi \leqslant \omega < \pi\) (r \textless{} 0, \(0 < \omega < \pi\) )。设 G 是由 u 生成的 X 的同胚所成的群，并赋 G 以离散拓扑。证明 G 在 X 上等度连续，且 X 中使 G 真作用的点所成之集是半平面 y \textgreater{} 0。
\end{enumerate}

\begin{enumerate}
\def\labelenumi{\arabic{enumi})}
\setcounter{enumi}{18}
\tightlist
\item
  设 X 是豪斯多夫一致空间，G 是 X 到其自身上的同胚所成的一个群。
\end{enumerate}

\begin{enumerate}
\def\labelenumi{\alph{enumi})}
\item
  证明若 \(\mathbf{G}\) 关于逐点收敛拓扑是等度连续且离散的，则 \(\mathbf{G}\) 在 \(\mathcal{F}_s(\mathbf{X};\mathbf{X})\) 中是闭的。
\item
  证明若 G（赋以离散拓扑）至少在 X 的一点处真作用（习题 18），则 G 在 \(\mathcal{F}_{s}(\mathbf{X};\mathbf{X})\) 中是闭的，且逐点收敛拓扑在 G 上诱导离散拓扑。
\item
  设 X 是两个空间 \(X_{1}, X_{2}\) 的拓扑和，它们每个都同胚于 R，从而 X 可等同于乘积空间 \(R \times \{1, 2\}\) 。赋 X 以乘积一致结构。对每对有理整数 \((m, n)\) ，设 \(u_{mn}\) 表示 X 的由 \(u_{mn}(x_{1}) = x_{1} + m\alpha + n\beta\) ， \(u_{mn}(x_{2}) = x_{2} + m\gamma + n\delta\) 定义的同胚，其中 \(x_{1} \in X_{1}\) 且 \(x_{2} \in X_{2}\) ，而 \(\alpha, \beta, \gamma, \delta\) 是四个非零实数，使 \(\alpha/\beta\) 与 \(\gamma/\delta\) 是互不相同的无理数。证明 \(u_{mn}\) 构成 X 的一个同胚群，它关于逐点收敛拓扑是等度连续且离散的，但 G（赋以离散拓扑）不在 X 的任何点处真作用。
\end{enumerate}

\begin{enumerate}
\def\labelenumi{\arabic{enumi})}
\setcounter{enumi}{19}
\tightlist
\item
  设 G 是豪斯多夫一致空间 X 的同胚所成的一个群。假设 G 是等度连续的，且 G（赋以离散拓扑）在 X 上真作用。再假设 \(x \in X\) 中不是异于恒等映射的任何同胚 \(u \in G\) 的不动点的那些点全体所成之集在 X 中稠密。在这些条件下，证明存在 X 中的一个开集 F，使得 \(F \cap U(F) = \varnothing\) 对异于恒等映射的每个 \(u \in G\) 成立，且 F 在轨道空间 X/G 中的典范象是 X/G 的一个稠密开子集，同胚于 F（利用习题 16 与佐恩引理）。
\end{enumerate}

